\section{Guarantees and Correctness Properties}\label{sec:properties}

The preceding sections establish the Bailout Protocol as a formal mechanism: a deterministic state machine $\mathcal{B} = (Q, q_0, \Sigma, \rightarrow, Q_F)$, a compulsory trigger condition, a parent-chain bypass mechanism, and an immutable escalation pathway to a named human principal. This section formally states the three correctness properties the protocol is designed to guarantee, identifies the structural assumptions on which each guarantee rests, and provides a proof sketch for each. All three properties are stated as invariants over the full protocol state space — they hold for every reachable state, not merely for a specific execution path.

% ---------------------------------------------------------------
\subsection{Property 1: Safety — No Autonomous Collapse}
% ---------------------------------------------------------------

\theorem[Safety (No Autonomous Collapse)]{During any execution of the Bailout Protocol, the system shall not transition to a terminal resolving state without an explicit, recorded human authorization.}

\subparagraph{Formal statement.} Let $\mathcal{B}$ be the Bailout Protocol state machine for node $n$. Let $P(t)$ denote the set of functional communication pathways to $h(n)$ available at time $t$. Safety requires that for all times $t$:

\begin{equation}
\forall t : |P(t)| = 0 \;\land\; \text{state}(n) = \text{RESOLVED}
\;\implies;
\text{tag}(n) = \texttt{unresolved}
\end{equation}

Equivalently: if no pathway to the human anchor is available, the system may enter \texttt{RESOLVED} only with an \texttt{unresolved} tag — a documented stall, not an autonomous resolution.

\subparagraph{Structural assumptions.}
\begin{enumerate}
\item[\textbf{A1.1}] The Fact Layer pre-commit hook is operative: any \texttt{UpdateLedger} call that would record a directive whose issuer is not the verified human principal $h(n)$ is rejected at the authorization layer.
\item[\textbf{A1.2}] The immutable parent-pointer $\alpha(n)$ and human-root identifier $h(n)$ are established at node provisioning and are not writable by any machine actor at runtime.
\item[\textbf{A1.3}] At least one functional communication pathway to $h(n)$ is provisioned at node startup. This is a deployment requirement, not a protocol invariant.
\end{enumerate}

\subparagraph{Proof sketch.} The proof proceeds by enumerating all transitions into the terminal state \texttt{RESOLVED}.

\begin{itemize}
\item[\textbf{Step 1.}] The only transitions entering \texttt{RESOLVED} are: T7 ($\text{ESCALATING} \rightarrow \text{RESOLVED}$ via human \texttt{RESOLVE} or \texttt{REJECT}) and T9 ($\text{HUMAN\_ACKNOWLEDGED} \rightarrow \text{RESOLVED}$ via human \texttt{RESOLVE} or \texttt{REJECT}). Both require a directive $d \in \{\texttt{RESOLVE}, \texttt{REJECT}\}$ issued by $h(n)$. By INV-2 (the Fact Layer enforcement invariant: $\text{state}(n) = \text{HUMAN\_ACKNOWLEDGED} \implies \text{directive\_issuer}(n) = h(n)$), the Fact Layer's pre-commit hook rejects any resolution attempt in which the directive's issuer field does not match $h(n)$. A machine actor cannot produce a directive that passes this gate.

\item[\textbf{Step 2.}] Consider the two source states for \texttt{RESOLVED}:
\begin{enumerate}
\item[\textbf{From ESCALATING:}] The \texttt{Escalate} algorithm propagates the exception along the parent chain using the Pathway Health Registry. If all provisioned pathways fail $N_{\max}$ consecutive health checks, the algorithm fires a \texttt{parent\_unreachable} timeout event, records the failure in the Forensic Ledger, and the originating node stalls in \texttt{ESCALATING}. No machine actor can issue a directive from this state by the bypass property (BP): the bypass compels all directives from \texttt{ESCALATING} to originate from $h(n)$. Therefore \texttt{RESOLVED} cannot be entered without human authorization when all pathways are unavailable.

\item[\textbf{From HUMAN\_ACKNOWLEDGED:}] Entry to \texttt{HUMAN\_ACKNOWLEDGED} occurs only via T2 or T6, both requiring a Fact-Layer-recorded human acknowledgment. Only $h(n)$ may issue a directive from this state. If the human anchor does not respond before $T_{\text{human}}$ expires, the protocol transitions to \texttt{RESOLVED} with an \texttt{unresolved} tag (T10) — a terminal state that documents a stall condition, not an autonomous resolution.
\end{enumerate}

\item[\textbf{Step 3.}] The trigger condition TC and the bypass mechanism BP jointly ensure that no transition in $Q \setminus \{\text{RESOLVED}\}$ can be reached via a machine-actor-issued directive. Reaching \texttt{RESOLVED} without authorization would require the Fact Layer to approve a directive from a non-human actor, which violates A1.1. Safety thus holds conditionally on A1.1.
\end{itemize}

\subparagraph{Discussion.} Safety does not require that a pathway to $h(n)$ always exists — it requires that the system cannot autonomously collapse regardless of pathway availability. When all pathways fail, the system stalls in \texttt{ESCALATING} or enters \texttt{RESOLVED} with \texttt{unresolved}. Both outcomes are documented stalls, not unauthorized resolutions. In safety-critical domains, a system that fails by stalling and logging the failure is preferable to one that proceeds on an unauthorized resolution. The safety property guarantees that the Bailout Protocol's failure mode is always the safer one.

% ---------------------------------------------------------------
\subsection{Property 2: Liveness — Eventual Human Contact}
% ---------------------------------------------------------------

\theorem[Liveness (Eventual Human Contact)]{During any execution of the Bailout Protocol, if the system enters the \texttt{BOUNDED} state and the Bounds Engine produces no Fact-Layer-approved resolution within $T_{\text{bounds}}$, then the protocol must eventually enter either \texttt{HUMAN\_ACKNOWLEDGED} or \texttt{RESOLVED} (with an \texttt{unresolved} tag), and in either case the human anchor $h(n)$ is contacted within a bounded wall-clock interval from the trigger event.}

\subparagraph{Formal statement.} Let $t_0$ be the time at which T1 fires (entry to \texttt{BOUNDED}). Liveness requires:

\begin{equation}
\exists t_1 \geq t_0 : t_1 - t_0 \leq T_{\max}
\;\implies\;
\text{state}(n) \in \{\text{HUMAN\_ACKNOWLEDGED}, \text{RESOLVED}\}
\end{equation}

where $T_{\max} \leq T_{\text{bounds}} + N_{\max} \cdot T_{\text{cap}} + T_{\text{human}}$ is a deployment-specific bound. \texttt{RESOLVED} with an \texttt{unresolved} tag counts as human contact: the notification was sent and received, as evidenced by the Forensic Ledger entry.

\subparagraph{Structural assumptions.}
\begin{enumerate}
\item[\textbf{A2.1}] The Pathway Health Registry is operative: at least one functional pathway to $h(n)$ is available at the time of the trigger event, or becomes available before $N_{\max}$ retries are exhausted.
\item[\textbf{A2.2}] The human anchor $h(n)$ is provisioned with a finite response timeout $T_{\text{human}}$. Non-response within $T_{\text{human}}$ triggers a reminder notification and transitions the protocol to \texttt{RESOLVED} with an \texttt{unresolved} tag, recorded in the Forensic Ledger.
\item[\textbf{A2.3}] The timeout parameters $T_{\text{bounds}}$, $T_{\text{prop}}$, $T_{\text{cap}}$, and $T_{\text{human}}$ are finite positive deployment constants, not adjustable by any machine actor at runtime.
\end{enumerate}

\subparagraph{Proof sketch.} The proof establishes that the protocol always advances from \texttt{BOUNDED} to a terminal state within bounded time, by examining each transition path and its timing constraints.

\begin{itemize}
\item[\textbf{Step 1.}] From \texttt{BOUNDED}, exactly two exits exist: T2 ($\text{BOUNDED} \rightarrow \text{HUMAN\_ACKNOWLEDGED}$ via intra-bounds resolution) or T3 ($\text{BOUNDED} \rightarrow \text{BAIL\_PENDING}$ upon $T_{\text{bounds}}$ expiration or explicit \texttt{bailout\_signal}). T2 is the fast path; T3 is the guaranteed path. Since $T_{\text{bounds}}$ is a hard finite timeout (A2.3), T3 fires within $T_{\text{bounds}}$ if T2 has not fired. The Fact Layer enforces $T_{\text{bounds}}$ at the pre-commit hook; no machine actor can prevent its expiration.

\item[\textbf{Step 2.}] T4 ($\text{BAIL\_PENDING} \rightarrow \text{ESCALATING}$) fires immediately upon entry to \texttt{BAIL\_PENDING} with no configurable dwell time — an architectural compulsion, not a configurable delay. From the moment T3 fires, the protocol enters \texttt{ESCALATING} within a bounded atomic transition latency. The Fact Layer pre-commit hook enforces this transition without machine-actor interposition.

\item[\textbf{Step 3.}] In \texttt{ESCALATING}, the \texttt{Escalate} algorithm iterates: it selects the lowest-latency functional pathway from the PHR, sends the exception to the parent node, and retries with exponential backoff capped at $T_{\text{cap}}$ if no acknowledgment arrives within $T_{\text{prop}}$. The maximum retry count $N_{\max}$ is a deployment constant. Total propagation time from \texttt{ESCALATING} entry to parent acknowledgment or $N_{\max}$-retry exhaustion is bounded by $N_{\max} \cdot T_{\text{cap}}$ plus accumulated backoff, which itself is bounded by $N_{\max} \cdot T_{\text{cap}}$.

\item[\textbf{Step 4.}] If the parent is unreachable after $N_{\max}$ retries at $T_{\text{cap}}$, the algorithm fires a \texttt{parent\_unreachable} timeout event, which propagates to $h(n)$ via the next available functional pathway in the PHR. This fallback pathway selection ensures that propagation does not deadlock on a single failed parent; it advances to $h(n)$ via an alternate route. Upon reaching $h(n)$, the protocol enters \texttt{HUMAN\_ACKNOWLEDGED} (T6).

\item[\textbf{Step 5.}] If $h(n)$ does not issue a directive before $T_{\text{human}}$ expires (A2.2), the protocol transitions to \texttt{RESOLVED} with an \texttt{unresolved} tag (T10). Both outcomes — human acknowledgment (\texttt{HUMAN\_ACKNOWLEDGED}) and human timeout with \texttt{unresolved} tag (\texttt{RESOLVED}) — constitute human contact as defined by the liveness property.
\end{itemize}

\subparagraph{Combined bound.} The total wall-clock time from $t_0$ (T1 firing) to human contact is bounded by:

\begin{equation}
T_{\max} \;\leq\; T_{\text{bounds}} \;+\; N_{\max} \cdot T_{\text{cap}} \;+\; T_{\text{human}}
\end{equation}

This bound is a deployment parameter, not a protocol constant. It depends on the network topology, the number of hops to $h(n)$, and the provisioned $T_{\text{human}}$. A deployment that cannot satisfy its operational $T_{\max}$ requirement must provision additional pathways, increase $N_{\max}$, or reduce timeout values — but the protocol guarantees human contact within whatever bound the deployment provisions.

\subparagraph{Discussion.} Liveness is a conditional guarantee: it holds if A2.1, A2.2, and A2.3 are simultaneously satisfied. If the human anchor is permanently unreachable — a deployment failure violating A2.1 — the protocol stalls in \texttt{ESCALATING} and records the failure. Safety (Property 1) ensures this stall does not become an autonomous resolution. Together, Safety and Liveness define the protocol's failure semantics: when human contact cannot be achieved because all pathways have failed, Safety guarantees the system fails by stalling, never by proceeding without authorization.

% ---------------------------------------------------------------
\subsection{Property 3: Accountability — Human Always Identifiable}
% ---------------------------------------------------------------

\theorem[Accountability (Human Always Identifiable)]{For every protocol run that produces a directive $d \in \{\texttt{ACK}, \texttt{RESOLVE}, \texttt{REJECT}\}$, the Forensic Ledger entry recording $d$ identifies the human principal who issued it, and no ledger entry records a directive issued by a machine actor.}

\subparagraph{Formal statement.} Let $\mathcal{L}$ be the append-only Forensic Ledger recording all protocol events. Accountability requires:

\begin{equation}
\forall \ell \in \mathcal{L} : \ell.\text{event\_type} = \text{human\_directive}
\;\implies;
\ell.\text{principal\_id} = h(n)
\;\land\;
\text{is\_human}(\ell.\text{principal\_id})
\end{equation}

where $h(n)$ is the human accountability anchor of the originating node $n$, and $\text{is\_human}$ returns true only for identifiers classified as human principal identifiers by the provisioning authority. Equivalently: every directive in the ledger is attributable to a human principal, and no machine-issued directive is recorded as valid.

\subparagraph{Structural assumptions.}
\begin{enumerate}
\item[\textbf{A3.1}] The provisioning authority issues human principal identifiers only to verified natural persons or institutional actors acting with unambiguous human authority. The $\text{is\_human}$ predicate is a deployment-side classification, not a protocol mechanism.
\item[\textbf{A3.2}] The Fact Layer pre-commit hook is operative: any \texttt{UpdateLedger} call that would record a directive with a non-human $\text{principal\_id}$ is rejected at the authorization layer.
\item[\textbf{A3.3}] The parent pointer $\alpha(n)$ and human-root identifier $h(n)$ are immutable for the node's lifetime. Any attempted modification at runtime is rejected by the Fact Layer pre-commit hook and logged as a protocol violation.
\end{enumerate}

\subparagraph{Proof sketch.} The proof proceeds by structural induction over the set of protocol transitions that can produce a ledger entry with $\text{event\_type} = \text{human\_directive}$, showing that only human-issued directives can produce a valid entry.

\begin{itemize}
\item[\textbf{Step 1.}] The only protocol transitions that generate a ledger entry with $\text{event\_type} = \text{human\_directive}$ are: T6 ($\text{ESCALATING} \rightarrow \text{HUMAN\_ACKNOWLEDGED}$ upon human \texttt{ACK}), T7 ($\text{ESCALATING} \rightarrow \text{RESOLVED}$ upon human \texttt{RESOLVE} or \texttt{REJECT}), T9 ($\text{HUMAN\_ACKNOWLEDGED} \rightarrow \text{RESOLVED}$ upon human \texttt{RESOLVE} or \texttt{REJECT}), and T10 ($\text{HUMAN\_ACKNOWLEDGED} \rightarrow \text{RESOLVED}$ upon human timeout with \texttt{unresolved} tag). Each of these transitions fires only upon receipt of a directive from $h(n)$, as enforced by the Fact Layer pre-commit hook and the authorization ledger.

\item[\textbf{Step 2.}] The Fact Layer pre-commit hook enforces INV-2: upon directive receipt, it verifies the principal identifier against the authorization ledger. If the identifier does not satisfy $\text{is\_human}$ or does not match the stored $h(n)$ for the originating node, the pre-commit hook rejects the directive and logs a protocol violation. This check cannot be bypassed by any machine actor, including the Purpose Layer: the Purpose Layer's role is exclusively receiving exceptions and issuing directives; it cannot authorize ledger entries attributing a directive to a non-human actor.

\item[\textbf{Step 3.}] The append-only property of $\mathcal{L}$ ensures that once written, no ledger entry can be modified or deleted. Even if a machine actor could produce a directive with a falsified principal identifier, the Fact Layer's pre-commit check prevents the entry from being written. Falsification is blocked at the authorization layer, not detected after the fact.

\item[\textbf{Step 4.}] The immutability of $h(n)$ (A3.3) ensures that the accountability anchor for each node is fixed at provisioning and cannot be redirected by a compromised or malfunctioning node at runtime. This prevents a machine actor from substituting its own human anchor for the correct one. The $h(n)$ field is stored in Fact Layer-managed memory and is read-only at the machine-actor level; any modification attempt is rejected and logged as an unauthorized modification.
\end{itemize}

\subparagraph{Discussion.} The accountability property guarantees attribution, not correctness: it guarantees that every directive in the ledger is attributable to a human principal, but not that the human principal made the correct decision. The latter is a separate concern — of human factors, organizational process, and training — that the protocol cannot address architecturally. What the protocol guarantees architecturally is that no machine actor can issue, simulate, or substitute a directive without this being detectable as a protocol violation. The ledger provides the documentary evidence; the pre-commit hook provides the enforcement. Together they make accountability non-repudiable: a human principal cannot later deny issuing a directive that appears in the ledger under their identifier, and a machine actor cannot cause a directive to appear under a falsified human identifier.

The three properties are structurally independent but jointly necessary. Safety prevents the system from entering a terminal state without human authorization. Liveness ensures the system reaches human contact within a bounded time if a pathway exists. Accountability ensures that when human contact occurs, the resulting directive is attributable to a human and not to any machine actor. Together they define the upstream accountability guarantee that the Bailout Protocol exists to enforce: every unresolved exception in a \bxthree{} system reaches a named human principal, by compulsion, within bounded time, with full forensic attribution — or the system stalls, safely, in a documented and auditable state.\footnote{This section formally derives properties referenced in \S\ref{sec:bailout} and \S\ref{sec:implementation}.}
